\documentclass[12pt]{report}

% ------------------------------------------------
% PACOTES
% ------------------------------------------------

\usepackage[utf8]{inputenc}
\usepackage[T1]{fontenc}
\usepackage[brazil]{babel}
\usepackage{amsmath}
\usepackage{amssymb}
\usepackage{graphicx}
\usepackage{booktabs}
\usepackage{hyperref}
\usepackage{float}
\usepackage{setspace}

\onehalfspacing

% ------------------------------------------------
% INFORMAÇÕES DO TRABALHO
% ------------------------------------------------

\title{Alocação Adaptativa de Portfólio Utilizando Aprendizado de Máquina e Otimização Média-Variância}
\author{Maria Eduarda Mesquita Magalhães}
\date{2026}

% ------------------------------------------------
% DOCUMENTO
% ------------------------------------------------

\begin{document}

\maketitle

\tableofcontents

\newpage

% =================================================
% 1 INTRODUÇÃO
% =================================================

\chapter{Introdução}

\section{Motivação}

A alocação de portfólio é um problema central em finanças quantitativas. Os investidores buscam distribuir seu capital entre diferentes ativos de forma a equilibrar retorno esperado e risco.

Ao longo das últimas décadas, a gestão de investimentos passou por transformações significativas impulsionadas pelo avanço dos métodos quantitativos e das técnicas computacionais. Em particular, os modelos de aprendizado de máquina têm sido cada vez mais utilizados nos mercados financeiros para melhorar previsões e apoiar processos de tomada de decisão.

Apesar desses avanços, muitas técnicas de construção de portfólio ainda se baseiam em estruturas clássicas, como a Teoria Moderna do Portfólio proposta por Harry Markowitz.

\section{Teoria Moderna do Portfólio}

A Teoria Moderna do Portfólio (Modern Portfolio Theory -- MPT), introduzida por Harry Markowitz, fornece uma estrutura matemática para construir portfólios que maximizem o retorno esperado para um determinado nível de risco.

O problema clássico de otimização média-variância é definido por:

\begin{equation}
\min_w w^T \Sigma w
\end{equation}

sujeito a:

\begin{equation}
w^T \mu = \mu_p
\end{equation}

onde:

\begin{itemize}
\item $w$ representa o vetor de pesos do portfólio;
\item $\Sigma$ é a matriz de covariância dos retornos dos ativos;
\item $\mu$ é o vetor de retornos esperados;
\item $\mu_p$ é o retorno alvo do portfólio.
\end{itemize}

\section{Problema de Pesquisa}

Embora a otimização média-variância seja teoricamente elegante, seu desempenho depende fortemente da qualidade das estimativas dos retornos esperados e das covariâncias.

Isso leva à questão central deste trabalho:

\textit{Modelos de aprendizado de máquina podem melhorar a alocação de portfólio em relação à otimização média-variância tradicional?}

\section{Contribuição}

Esta dissertação propõe uma estrutura híbrida que combina modelos de aprendizado de máquina com técnicas clássicas de otimização de portfólio.

O pipeline proposto consiste em:

\begin{enumerate}
\item Utilizar modelos de aprendizado de máquina para prever retornos esperados dos ativos;
\item Estimar a matriz de covariância dos retornos;
\item Aplicar a otimização média-variância para determinar os pesos do portfólio;
\item Comparar o desempenho com métodos tradicionais de alocação.
\end{enumerate}

\newpage

% =================================================
% 2 REVISÃO DA LITERATURA
% =================================================

\chapter{Revisão da Literatura}

\section{Teoria Clássica de Portfólios}

A otimização média-variância representa uma das estruturas mais influentes da economia financeira. Ela formaliza a relação entre risco e retorno por meio do conceito de fronteira eficiente.

A fronteira eficiente representa o conjunto de portfólios que alcançam o maior retorno esperado para um determinado nível de risco.

\section{Índice de Sharpe e CAPM}

William F. Sharpe introduziu o Índice de Sharpe como uma medida de retorno ajustado ao risco:

\begin{equation}
S = \frac{R_p - R_f}{\sigma_p}
\end{equation}

onde:

\begin{itemize}
\item $R_p$ representa o retorno do portfólio;
\item $R_f$ representa a taxa livre de risco;
\item $\sigma_p$ representa a volatilidade do portfólio.
\end{itemize}

O Modelo de Precificação de Ativos Financeiros (CAPM) estende essas ideias ao relacionar retornos esperados ao risco sistemático de mercado.

\section{Limitações da Otimização Média-Variância}

Apesar de sua elegância teórica, a abordagem de Markowitz apresenta diversas limitações.

\subsection{Erro de Estimação}

Os retornos esperados são notoriamente difíceis de estimar. Pequenos erros podem gerar grandes mudanças nos pesos ótimos dos ativos.

\subsection{Sobreajuste (Overfitting)}

A otimização pode amplificar ruídos estatísticos presentes nos dados, produzindo alocações extremas e pouco robustas.

\subsection{Instabilidade das Matrizes de Covariância}

As correlações entre ativos variam ao longo do tempo, tornando a estimação da matriz de covariância instável.

\subsection{Distribuições Não Normais}

Os retornos financeiros frequentemente apresentam assimetria e caudas pesadas, violando a hipótese de normalidade assumida em muitos modelos clássicos.

\newpage

% =================================================
% 3 APRENDIZADO DE MÁQUINA EM FINANÇAS
% =================================================

\chapter{Aprendizado de Máquina em Finanças}

As técnicas de aprendizado de máquina têm recebido crescente atenção em aplicações financeiras, incluindo previsão de retornos, modelagem de risco e seleção de ativos.

\section{Aprendizado Supervisionado}

No aprendizado supervisionado, os modelos são treinados utilizando dados históricos compostos por variáveis explicativas e variáveis-alvo.

Fluxo típico:

\begin{itemize}
\item Conjunto de treinamento;
\item Conjunto de validação;
\item Conjunto de teste.
\end{itemize}

\section{Modelos}

Diversos modelos podem ser utilizados para tarefas de previsão financeira.

\subsection{Modelos Lineares}

\begin{itemize}
\item Regressão Linear;
\item Regressão Ridge;
\item Regressão Lasso.
\end{itemize}

\subsection{Modelos Não Lineares}

\begin{itemize}
\item Random Forest;
\item Gradient Boosting.
\end{itemize}

\subsection{Aprendizado Profundo}

\begin{itemize}
\item Redes Neurais;
\item Redes Long Short-Term Memory (LSTM).
\end{itemize}

\newpage

% =================================================
% 4 METODOLOGIA
% =================================================

\chapter{Metodologia}

\section{Dados}

A base de dados consiste em informações históricas de um conjunto de ativos financeiros.

Possível universo de ativos:

\begin{itemize}
\item Ações do mercado brasileiro (IBOVESPA);
\item ETFs globais.
\end{itemize}

As fontes de dados incluem:

\begin{itemize}
\item Yahoo Finance;
\item Kaggle;
\item Quandl.
\end{itemize}

O período amostral considerado neste estudo compreende os anos de 2010 a 2024.

\section{Engenharia de Atributos}

Diversas variáveis financeiras são construídas para alimentar os modelos de aprendizado de máquina.

Exemplos:

\begin{itemize}
\item Retornos passados;
\item Momentum de 1, 3 e 12 meses;
\item Volatilidade móvel;
\item Indicadores técnicos (RSI e médias móveis);
\item Variáveis macroeconômicas.
\end{itemize}

\section{Modelos de Aprendizado de Máquina}

O objetivo dos modelos é prever retornos futuros esperados:

\begin{equation}
E[r_{i,t+1}]
\end{equation}

Os seguintes modelos são avaliados:

\begin{itemize}
\item Regressão Linear;
\item Random Forest;
\item Gradient Boosting.
\end{itemize}

\section{Construção do Portfólio}

O processo de construção do portfólio consiste em três etapas principais:

\begin{enumerate}
\item Previsão dos retornos esperados por meio de aprendizado de máquina;
\item Estimação da matriz de covariância;
\item Resolução do problema de otimização de Markowitz.
\end{enumerate}

Uma formulação possível consiste em maximizar o Índice de Sharpe:

\begin{equation}
\max_w \frac{w^T \hat{\mu}}
{\sqrt{w^T \Sigma w}}
\end{equation}

\section{Estratégia de Backtesting}

Uma abordagem de janela móvel (rolling window) é utilizada para avaliação fora da amostra.

Exemplo:

\begin{itemize}
\item Período de treinamento: 2010--2016;
\item Período de teste: 2017.
\end{itemize}

A janela é então deslocada ao longo do tempo.

\newpage

% =================================================
% 5 AVALIAÇÃO DE DESEMPENHO
% =================================================

\chapter{Avaliação de Desempenho}

A estratégia proposta é comparada com diferentes portfólios de referência.

\section{Benchmarks}

\begin{itemize}
\item Portfólio igualmente ponderado;
\item Portfólio clássico de Markowitz;
\item Portfólio Aprendizado de Máquina + Markowitz.
\end{itemize}

\section{Métricas de Avaliação}

\subsection{Índice de Sharpe}

\begin{equation}
S = \frac{R_p - R_f}{\sigma_p}
\end{equation}

\subsection{Máximo Drawdown}

O máximo drawdown mede a maior perda acumulada entre um pico e um vale durante o período analisado.

\subsection{Retorno Acumulado}

Representa o crescimento total do portfólio ao longo do horizonte de investimento.

\subsection{Volatilidade}

Medida pelo desvio-padrão dos retornos do portfólio.

\newpage

% =================================================
% 6 RESULTADOS EMPÍRICOS
% =================================================

\chapter{Resultados Empíricos}

Esta seção apresenta os resultados obtidos nos experimentos de backtesting.

\section{Fronteira Eficiente}

Inserir gráfico da fronteira eficiente.

\begin{figure}[H]
\centering
\includegraphics[width=0.7\textwidth]{efficient_frontier.png}
\caption{Fronteira Eficiente}
\end{figure}

\section{Retornos Acumulados}

Inserir gráfico de retornos acumulados.

\section{Tabela de Desempenho}

\begin{table}[H]
\centering
\begin{tabular}{lcc}
\toprule
Modelo & Índice de Sharpe & Volatilidade \\
\midrule
Pesos Iguais & 0.70 & 0.18 \\
Markowitz & 0.80 & 0.16 \\
Aprendizado de Máquina + Markowitz & 1.10 & 0.15 \\
\bottomrule
\end{tabular}
\caption{Comparação de desempenho}
\end{table}

\newpage

% =================================================
% 7 TESTES DE ROBUSTEZ
% =================================================

\chapter{Testes de Robustez}

Experimentos adicionais são realizados para avaliar a robustez dos resultados.

\begin{itemize}
\item Diferentes períodos de análise;
\item Diferentes universos de ativos;
\item Inclusão de custos de transação.
\end{itemize}

\newpage

% =================================================
% 8 DISCUSSÃO
% =================================================

\chapter{Discussão}

Este capítulo interpreta os resultados empíricos obtidos.

Questões abordadas:

\begin{itemize}
\item Os modelos de aprendizado de máquina melhoram a previsão de retornos?
\item Os portfólios tornam-se mais estáveis?
\item O risco é reduzido?
\end{itemize}

\newpage

% =================================================
% 9 CONCLUSÃO
% =================================================

\chapter{Conclusão}

Esta dissertação investigou se modelos de aprendizado de máquina podem melhorar a alocação de portfólio quando combinados com a otimização média-variância.

Os resultados sugerem que abordagens híbridas que unem modelos preditivos e técnicas clássicas de otimização podem melhorar o retorno ajustado ao risco.

\section{Limitações}

\begin{itemize}
\item Limitações dos dados;
\item Hipóteses dos modelos;
\item Custos de transação.
\end{itemize}

\section{Pesquisas Futuras}

Possíveis extensões incluem:

\begin{itemize}
\item Modelos de aprendizado de máquina para estimação de covariâncias;
\item Aprendizado por reforço para alocação de portfólio;
\item Arquiteturas avançadas de deep learning.
\end{itemize}

\end{document}